\chapter{Cosmology and the Limits of the Frame}
\label{chap:cosmology-limits-of-the-frame}

Part V closes at the limits of the frame. Three effects bring it there: the dark night sky that a finite count
explains, the flat rotation curves that a missing count only names, and --- the deepest no-go in the book ---
the closed loop in time that the causal order forbids. The chapter ends where the frame itself ends: a causal
order cannot close on itself, because the count can only rise.

\section{The Olbers Effect: a finite count makes a dark sky}
\label{se:olbers}

Why is the night sky dark? The question sounds trivial and is not. In an infinite, eternal, static universe
the sky should blaze: look in any direction and your line of sight, continued far enough, must eventually
strike the surface of some star, so every point of the sky should shine as bright as a stellar surface, and
night should be as bright as day. The naive continuum makes the argument precise. Counting the light from
shells of sources at every distance, the brightness integrates as
\begin{equation}
  \int (\text{source density})\, dr,
\end{equation}
and over an infinite static universe filled uniformly with stars the integral diverges. The cancellation is the
heart of it: a shell of sources at radius $r$ has an area that grows as $r^2$, so it holds about $r^2$ times as
many stars as a nearby shell, while each of those stars is dimmed by the inverse square of its distance, a
factor of $1/r^2$. The two run exactly opposite and cancel, so every shell, near or far, delivers the same
total light to the eye --- and an infinity of equal contributions sums without bound, predicting a sky uniformly
ablaze.

The night sky is dark, and the resolution is a finite count. The paradox needed three things at once --- a
universe infinite, eternal, and static --- and the real one is none of them. It has a finite age: the causal
record reaches only as far as light has had time to travel since the beginning, a horizon of size $c$ times the
age of the universe, about fourteen billion light-years, and within that horizon there are only finitely many
sources. The shells do not run on forever; they stop at the horizon, and a finite sum of finite contributions
is finite, and dark. And the universe is not static but expanding, which thins the count a second way: the
light of the most distant sources is redshifted as it crosses the growing space, sliding down out of the
visible and dimming further, so that even the sources inside the horizon do not all contribute their full
share. Both corrections pull the same direction --- toward fewer effective sources, a count made finite. Nor
does intervening dust save the static picture: dust set between the stars would absorb their light, warm, and
in an eternal universe come to glow as bright as what it blocked, so absorption only redistributes the blaze;
only the finiteness of the count removes it. Put thermodynamically the paradox is sharper still: in an eternal
universe the starlight would long since have come into equilibrium, filling all of space with radiation at the
temperature of a stellar surface --- thousands of degrees --- the night sky not merely bright but hot. The sky
is instead dark and cold, a few degrees above absolute zero, and that coldness is the same finite count read as
a temperature: too few sources, shining for too short a time, to fill the sky with their heat. Edgar Allan Poe
guessed as much in 1848, writing that the gaps
between the stars must be places whose light had not yet had time to reach us --- the finite-age resolution,
intuited by a poet before the cosmologists held it.

The honest floor is a finite count obligation --- the night sky is dark because the count of sources within the
causal horizon is finite. The reading is that Olbers' paradox dissolves the moment the count is finite. The
dark sky is not a puzzle but a measurement: it shows, directly, that the universe a record can see holds
finitely many sources, the darkness between the stars the empty slots a finite count leaves. An infinitely
refined ledger would admit infinitely many luminous events; the darkness is the direct observational proof that
the record is finite.

\section{The Flat Rotation Effect: a missing count, named}
\label{se:flat-rotation}

Stars at the edge of a galaxy orbit too fast. For a star at radius $r$, the gravity holding it in orbit comes
from the mass enclosed within that radius, and balancing the pull against the orbital motion gives
\begin{equation}
  v^2(r) = \frac{G M(r)}{r},
\end{equation}
so once you are past the bright central bulge, where the enclosed mass $M(r)$ stops growing, the speed should
fall off as
\begin{equation}
  v(r) \propto r^{-1/2},
\end{equation}
the Keplerian decline the planets obey, each outer one slower than the last. The measured curves do not fall.
They stay flat,
\begin{equation}
  v(r) \to \text{const},
\end{equation}
far out where the visible mass has run out --- and a flat $v$ with $v^2 = GM(r)/r$ forces $M(r)$ to keep
growing in proportion to $r$, mass piling up linearly with radius long after the light has given out. Something
unseen, a halo of it, holds the fast stars in. Vera Rubin and her collaborators measured these flat curves,
galaxy by galaxy, through the 1970s, and the discrepancy between the flat curve and the Keplerian one is
conventionally recorded as ``dark matter.'' The flat curves were not the first sign of it: Fritz Zwicky, in
1933, found the galaxies of the Coma cluster swarming far too fast for the visible mass to bind them, and named
the shortfall --- \emph{dunkle Materie}, dark matter --- four decades before the rotation curves confirmed it
one galaxy at a time.

The gap admits more than one reading, and the model takes no side. One reading keeps the law of gravity and adds
the mass: a halo of unseen matter, gravitating but not shining, outweighing the stars several times over. The
other keeps the visible mass and changes the law: gravity departs from the inverse square at the very small
accelerations found at a galaxy's rim, so that no missing matter is needed. The weight of the wider evidence ---
the same unseen mass bending light into gravitational lenses, and, in the Bullet Cluster, two clusters caught
just after a collision with the lensing mass found sailing ahead of the visible gas that collided and lagged
--- has moved most astronomers toward the first reading; but that adjudication is physics, and the finite model
makes neither call.

The honest floor is name-only --- the missing count is named, not derived. The finite model records the
discrepancy and gives it a name; it does not derive the missing mass, nor settle whether ``dark matter'' is the
right name for it, nor whether the fault lies in the unseen mass or in the law of gravity itself. The reading
is that the flat rotation curve names a missing count without explaining it. The stars orbit as if more mass
were present than is seen, the curve stays flat where it should fall, and the model does the honest minimum: it
records the gap and labels it, leaving to physics what the label stands for.

\section{The Time Travel Effect: the frame cannot close on itself}
\label{se:time-travel}

The book's deepest no-go closes the part, and it is deeper than the others because it is not a wall against
some particular refinement but a wall against the frame closing on itself. A causal order is acyclic. Each
event follows the events it refines and precedes the events that refine it, and the relation is a strict
partial order: irreflexive, so nothing precedes itself; transitive, so the chain of precedence composes; and
asymmetric, so if one event precedes another the second cannot precede the first. A record cannot be its own
predecessor any more than a line of a ledger can be entered before itself. A closed timelike curve --- a loop
in time, a path that returns to its own past --- would require exactly that, an event preceding itself, and the
order forbids it outright. This is the grandfather paradox stated structurally: the loop is impossible not for
what it would let you do once you arrived, but because the order has no room for a cycle.

Try to force the loop and the bookkeeping breaks. Bending a finite causal chain back into itself,
\begin{equation}
  U_{n+1} = U_n + \int_0^1 E_{j,\alpha}\, d\alpha,
\end{equation}
demands an uncountable set of events inserted between an event and itself, destroying the local finiteness
every record must have --- a finite chain cannot be closed into a circle without an infinity of new links where
there is room for none. There is no admissible refinement that closes the loop. One might hope to rescue the loop by asking only for
self-consistent histories --- that whatever returns to the past was always part of it, so that nothing is ever
undone --- but consistency is not where the wall stands. The wall is prior to it: a strict order has no cycle
for consistency to be imposed on, and the count that fixes the order only ever rises.

The book has met many walls --- the slip below the friction threshold, the transmission through crossed
polarizers, the seam through matched domains, the path through Anderson's disorder, the factorization of Bell's
correlations, the outbound record across an event horizon --- each a particular refinement the law forbids.
This last is their root. Those were walls within the frame, against one refinement or another; this is a wall
against the frame itself, against the order that underwrites them all. It is the deepest because it is the one
the others rest on: a causal order that could loop could forbid nothing, for in a cycle every event would both
precede and follow every other, and the very notion of an admissible refinement would dissolve.

And the reason beneath the reason is the arrow. The count can only rise, as the arrow of time in this part's
first chapter required: each tick closes a degree of freedom, and a closed degree cannot reopen, so the count
grows and never returns to a value it has left. A closed timelike curve would demand exactly that return --- the
count coming back to where it was --- and a count that only rises never does. The honest floor is a no-go, and
it is the frame's own deepest one. The reading is that time cannot close on itself. The same fact that gave time
its arrow forbids the loop: a circle in time would need the count to come back to where it was, and a count that
only rises never does. The deepest wall in the book is the one against a circle in time.

\bigskip

With the frame's refusal to close, Part V ends. Relativity and gravitation, read off finite records, are the
geometry of the count: the interval the count agrees on across frames, the potential a local rescaling of the
count rate, the horizon a wall the count cannot cross, and the closed loop in time a no-go the count's own
rising forbids. The geometry that bends in this part is the geometry of counting itself --- not a stage the
count moves through but the shape the count makes. The last part turns inward, away from any particular physics
and toward the foundations: what the apparatus, whatever it measures, cannot escape --- the floor beneath every
floor, the assumptions no measurement can do without.
